\documentclass[11pt,a4paper]{article}
\usepackage[utf8]{inputenc}
\usepackage[spanish]{babel}
\usepackage{geometry}
\usepackage{booktabs}
\usepackage{tabularx}
\usepackage{xcolor}
\usepackage{hyperref}

\geometry{top=2.5cm, bottom=2.5cm, left=2.5cm, right=2.5cm}
\definecolor{brandgreen}{RGB}{27, 77, 62}

\title{\textbf{\color{brandgreen} PDAEH -- Plantilla de Artefacto: DPV} \\[0.3em] \large Documento de perspectivas y vistas arquitectónicas}
\author{Proceso para la Definición de Arquitecturas Empresariales Híbridas}
\date{\today}

\begin{document}
\maketitle

\noindent \textit{\textbf{Nota sobre el uso de esta plantilla:} Esta plantilla es una guía metodológica sugerida y no constituye un instrumento rígido o mandatorio. Cada organización, empresa, grupo de investigación o equipo de trabajo es libre de adaptarla, extenderla o simplificarla según sus políticas internas, herramientas y nivel de madurez.}

\section*{1. Propósito y Alcance del Artefacto}
Definir las vistas y perspectivas que permiten representar el sistema ante diferentes tipos de interesados.

\section*{2. Estructura y Registro del Contenido}
\begin{tabularx}{\textwidth}{|>{\bfseries}p{4.5cm}|X|}
    \hline
    \rowcolor[gray]{0.92} Campo / Sección & Registro del Contenido \\ \hline
    \textbf{Perspectiva o vista} & [Completar información] \\ \hline
    \textbf{Propósito} & [Completar información] \\ \hline
    \textbf{Stakeholders} & [Completar información] \\ \hline
    \textbf{Preocupaciones atendidas} & [Completar información] \\ \hline
    \textbf{Contenido representado} & [Completar información] \\ \hline
    \textbf{Notación} & [Completar información] \\ \hline
    \textbf{Relaciones con otras vistas} & [Completar información] \\ \hline
    \textbf{Criterio de aceptación} & [Completar información] \\ \hline
\end{tabularx}

\vspace{1.2cm}
\section*{3. Control de Revisiones y Gobierno Arquitectónico}
\noindent \textbf{Elaborado por:} \hrulefill \\[0.5em]
\textbf{Fecha:} \hrulefill \\[0.5em]
\textbf{Aprobación del Órgano de Gobierno Arquitectónico (OGA):} Aprobado [ \, ] \quad Rechazado [ \, ]

\end{document}
